% !TeX program = xelatex
% Compile twice with XeLaTeX. All illustrations are native vector TikZ drawings.
\documentclass[12pt,a4paper,oneside]{report}
\XeTeXgenerateactualtext=1
\usepackage[margin=23mm,top=24mm,bottom=24mm,headheight=27pt]{geometry}
\usepackage{amsmath,amssymb,mathtools,bm}
\usepackage{booktabs,array,longtable,tabularx}
\usepackage[table]{xcolor}
\usepackage{tikz}
\usetikzlibrary{arrows.meta,calc,patterns}
\usepackage[most]{tcolorbox}
\usepackage{titlesec,needspace}
\usepackage{hyperref}
\usepackage{bookmark}
\usepackage{xepersian}
\settextfont[Scale=1.06,BoldFont=Amiri-Bold.ttf]{Amiri-Regular.ttf}
\setlatintextfont[BoldFont=texgyretermes-bold.otf,ItalicFont=texgyretermes-italic.otf,BoldItalicFont=texgyretermes-bolditalic.otf]{texgyretermes-regular.otf}
\definecolor{navy}{HTML}{163747}
\definecolor{teal}{HTML}{127C80}
\definecolor{gold}{HTML}{B18542}
\definecolor{paper}{HTML}{F3F7F7}
\definecolor{muted}{HTML}{5E7178}
\definecolor{alert}{HTML}{A45538}
\hypersetup{unicode=true,colorlinks=true,linkcolor=navy,urlcolor=teal,
 pdftitle={روش‌های آماری ـ جزوهٔ بازنویسی‌شده},pdfauthor={علیرضا جهانبخش و احسان شهبازی},pdfsubject={احتمال، برآورد نقطه‌ای، فاصلهٔ اطمینان و آزمون فرض}}
\setlength{\parindent}{0pt}
\setlength{\parskip}{5pt plus 1pt}
\linespread{1.12}
\raggedbottom
\setcounter{tocdepth}{1}
\setcounter{secnumdepth}{2}
\allowdisplaybreaks[2]
\emergencystretch=1em
% Page styles use explicit boxes to keep bidi direction groups balanced.
\makeatletter
\def\ps@lecture{%
  \def\@oddhead{\vbox{\hbox to\textwidth{\small
    \rl{روش‌های آماری}\hfil\rl{\leftmark}}%
    \kern4pt\hrule height .35pt}}%
  \let\@evenhead\@oddhead
  \def\@oddfoot{\hfil\thepage\hfil}%
  \let\@evenfoot\@oddfoot}
\makeatother
\pagestyle{lecture}
\renewcommand{\chaptermark}[1]{\markboth{#1}{}}
\titleformat{\chapter}[display]{\color{navy}\bfseries\Huge}{\large فصل \thechapter}{8pt}{}[\vspace{4pt}\color{gold}\titlerule]
\titlespacing*{\chapter}{0pt}{0pt}{18pt}
\titleformat{\section}{\Large\bfseries\color{navy}}{\thesection}{.6em}{}
\titleformat{\subsection}{\large\bfseries\color{teal}}{\thesubsection}{.6em}{}
\tcbset{enhanced,breakable,boxrule=.45pt,arc=1.5mm,
 left=3mm,right=3mm,top=2mm,bottom=2mm,before skip=9pt,after skip=9pt,
 fonttitle=\bfseries,coltitle=white,colback=paper,colframe=teal}
\newtcolorbox{definition}[1]{breakable=false,title={تعریف | #1},colframe=navy,colback=navy!3!white}
\newtcolorbox{result}[1]{breakable=false,title={نتیجهٔ کلیدی | #1},colframe=teal,colback=teal!4!white}
\newtcolorbox{important}[1]{breakable=false,title={دقت کنید | #1},colframe=gold,colback=gold!6!white}
\newcounter{ex}[chapter]
\newenvironment{example}[1]{\refstepcounter{ex}\begin{tcolorbox}[breakable=false,title={مثال \thechapter.\arabic{ex} | #1},colframe=muted,colback=white]}{\end{tcolorbox}}
\newcommand{\edit}[1]{\par\begingroup\small\color{alert}\textbf{تکمیل یا اصلاح علمی.} #1\par\endgroup}
\newcommand{\source}[1]{\par{\small\color{muted} مبنا: صفحه‌های #1 فایل دست‌نویس.}\par}
\newcommand{\E}{\mathbb E}
\newcommand{\Var}{\operatorname{Var}}
\newcommand{\Cov}{\operatorname{Cov}}
\newcommand{\MSE}{\operatorname{MSE}}
\newcommand{\Bias}{\operatorname{Bias}}
\newcommand{\SE}{\operatorname{SE}}
\newcommand{\iid}{\overset{\mathrm{iid}}{\sim}}
\newcommand{\N}{\mathcal N}
\newcommand{\R}{\mathbb R}
\newcommand{\Prob}{\mathbb P}
\newcommand{\sol}{\par\textbf{حل.} }
\newcommand{\figcap}[1]{\par\smallskip{\small\color{muted}#1}\par}

\hypersetup{pdftitle={آزمون فرض آماری},pdfauthor={دکتر سکینه دهقان؛ علیرضا جهانبخش و احسان شهبازی}}
\begin{document}
\setcounter{chapter}{3}
\chapter{آزمون فرض آماری}
\begin{center}{\small\color{muted} استاد درس: دکتر سکینه دهقان\quad | \quad گردآورندگان جزوه: علیرضا جهانبخش و احسان شهبازی}\end{center}

\source{۲۷ تا ۴۵}
\section{فرض آماری، آماره و ناحیهٔ رد}
\begin{definition}{اجزای یک مسئلهٔ آزمون فرض}
\textbf{فرض آماری} گزاره‌ای دربارهٔ توزیع جامعه یا پارامترهای آن است. فرض مبنا را با $H_0$ و فرض مقابل را با $H_1$ نشان می‌دهیم.

\textbf{فرض ساده} یک توزیع را به‌طور کامل مشخص می‌کند. فرضی که بیش از یک توزیع را مجاز بداند، مرکب است؛ حتی تثبیت یک پارامتر در حضور پارامتر مزاحم نامعلوم می‌تواند فرضی مرکب بسازد.

\textbf{آمارهٔ آزمون} تابعی از نمونه است که بر اساس مقدار آن دربارهٔ رد یا رد نکردن $H_0$ تصمیم می‌گیریم.

\textbf{ناحیهٔ رد} یا بحرانی، مجموعهٔ مقادیر آماره یا نمونه است که به رد $H_0$ منجر می‌شود؛ آن را با $C$ نشان می‌دهیم.
\end{definition}
برای مثال، در آزمون میانگین، سه فرض مقابل رایج عبارت‌اند از
\[
H_1:\mu<\mu_0,\qquad H_1:\mu>\mu_0,\qquad H_1:\mu\ne\mu_0.
\]
این‌ها به‌ترتیب آزمون چپ‌طرفه، راست‌طرفه و دوطرفه را تعیین می‌کنند. در هر سه حالت، مرز فرض صفر $\mu=\mu_0$ است. جهت فرض مقابل باید بر اساس پرسش علمی و پیش از دیدن داده‌ها تعیین شود.
\begin{important}{رد نکردن، اثبات نیست}
تصمیم معمول آزمون «رد $H_0$» یا «رد نکردن $H_0$» است. رد نکردن ممکن است از کم بودن شواهد یا توان پایین ناشی شود و به معنای اثبات $H_0$ نیست. در این نسخه، عبارت «قبول فرض صفر» در متن مبنا با همین تعبیر دقیق‌تر جایگزین شده است.
\end{important}
\section{خطاهای نوع اول و دوم؛ توان آزمون}
\source{۲۷ و ۲۸}
\begin{center}
\begin{tabular}{r|cc}
\toprule
تصمیم & $H_0$ درست است & $H_0$ نادرست است\\
\midrule
رد $H_0$ & خطای نوع اول & تصمیم درست\\
رد نکردن $H_0$ & تصمیم درست & خطای نوع دوم\\
\bottomrule
\end{tabular}
\end{center}
\begin{result}{سطح و توان}
برای یک فرض صفر ساده،
\[
\alpha=\Prob_{H_0}(\text{رد }H_0).
\]
برای یک مقدار مشخص $\theta\in\Theta_1$ از فرض مقابل،
\[
\beta(\theta)=\Prob_\theta(\text{رد نکردن }H_0),\qquad
\pi(\theta)=1-\beta(\theta)=\Prob_\theta(\text{رد }H_0).
\]
تابع $\pi(\theta)$ تابع توان است. برای فرض صفر مرکب، سطح آزمون با بیشترین احتمال خطای نوع اول روی مقادیر مجاز فرض صفر کنترل می‌شود.
\end{result}
معمولاً $\alpha$ را پیشاپیش در سطحی مانند $0.05$ تثبیت می‌کنیم و میان آزمون‌های دارای این کنترل، توان بیشتر را مطلوب می‌دانیم. با حجم نمونهٔ ثابت، کاهش $\alpha$ اغلب با افزایش $\beta$ همراه است؛ بزرگ‌تر کردن نمونه می‌تواند هر دو خطا را بهبود دهد.
\begin{center}
\begin{LTR}
\begin{tikzpicture}[x=1.2cm,y=4.1cm]
\fill[alert!20] plot[domain=1.645:4,samples=70] (\x,{exp(-\x*\x/2)/sqrt(2*pi)})--(4,0)--(1.645,0)--cycle;
\fill[teal!17] plot[domain=-1:1.645,samples=70] (\x,{exp(-(\x-2.5)^2/2)/sqrt(2*pi)})--(1.645,0)--(-1,0)--cycle;
\draw[navy,thick] plot[domain=-3:5.5,samples=120] (\x,{exp(-\x*\x/2)/sqrt(2*pi)});
\draw[teal,thick] plot[domain=-3:5.5,samples=120] (\x,{exp(-(\x-2.5)^2/2)/sqrt(2*pi)});
\draw[muted,-{Latex}] (-3.1,0)--(5.8,0);
\draw[gold,dashed] (1.645,0)--(1.645,.44);
\node[navy,above] at (0,.4) {$H_0$};\node[teal,above] at (2.5,.4) {$H_1$};
\node[below] at (1.645,0) {$c$};
\node[alert] at (2.15,.035) {$\alpha$};\node[teal] at (.95,.06) {$\beta$};
\end{tikzpicture}
\end{LTR}
\figcap{نمایش خطای نوع اول و دوم برای یک آزمون راست‌طرفه در برابر یک مقدار مشخص از فرض مقابل. ناحیهٔ رد، سمت راست $c$ است.}
\end{center}
\begin{example}{ده پرتاب سکه}
فرض کنید $X_i$ شاخص شیر آمدن در پرتاب $i$ باشد و
\[
X_i\iid\operatorname{Bern}(p),\quad n=10,\qquad
H_0:p=\frac12,\quad H_1:p=\frac34.
\]
ناحیهٔ رد را $C=\{\sum_iX_i\ge7\}$ انتخاب می‌کنیم. اگر $Y=\sum_iX_i$، آنگاه
\begin{align*}
\alpha&=\Prob_{p=1/2}(Y\ge7)
=\sum_{y=7}^{10}\binom{10}{y}2^{-10}
=\boxed{0.171875},\\
\beta&=\Prob_{p=3/4}(Y\le6)
=\sum_{y=0}^{6}\binom{10}{y}(3/4)^y(1/4)^{10-y}
\approx\boxed{0.224125}.
\end{align*}
توان در $p=3/4$ برابر $1-\beta\approx0.775875$ است. این آزمونِ ازپیش‌تعریف‌شده، سطح ۵ درصد ندارد؛ سطح واقعی آن همان $0.171875$ است.
\end{example}
\section{آزمون میانگین یک جامعهٔ نرمال؛ واریانس معلوم}
\source{۲۹ تا ۳۱}
فرض کنید $X_i\iid\N(\mu,\sigma^2)$ و $\sigma$ معلوم باشد. برای $H_0:\mu=\mu_0$، آمارهٔ
\[Z=\frac{\bar X-\mu_0}{\sigma/\sqrt n}\]
تحت فرض صفر، نرمال استاندارد است.
\subsection{استخراج ناحیهٔ رد چپ‌طرفه}
برای $H_1:\mu<\mu_0$، میانگین‌های کوچک به سود فرض مقابل‌اند. ناحیهٔ رد را $\bar X-\mu_0<c$ می‌گیریم. از
\[
\alpha=\Prob_{H_0}(\bar X-\mu_0<c)
=\Prob\left(Z<\frac{c}{\sigma/\sqrt n}\right)
\]
نتیجه می‌شود $c=z_\alpha\sigma/\sqrt n=-z_{1-\alpha}\sigma/\sqrt n$. پس ناحیهٔ رد $Z<-z_{1-\alpha}$ است.
\begin{result}{سه آزمون میانگین با واریانس معلوم؛ سطح $\alpha$}
با مقدار مشاهده‌شدهٔ $z_{\rm obs}$،
\[
\begin{array}{c|c}
\text{فرض مقابل}&\text{شرط رد فرض صفر}\\\hline
\mu<\mu_0&z_{\rm obs}<-z_{1-\alpha}\\[2pt]
\mu>\mu_0&z_{\rm obs}>z_{1-\alpha}\\[2pt]
\mu\ne\mu_0&|z_{\rm obs}|>z_{1-\alpha/2}
\end{array}
\]
\end{result}
برای آزمون راست‌طرفه، از $\Prob(Z<c/(\sigma/\sqrt n))=1-\alpha$ نتیجه می‌شود $c=z_{1-\alpha}\sigma/\sqrt n$. برای آزمون دوطرفه،
\[
\alpha=\Prob(|Z|>c/(\sigma/\sqrt n))
=2\left[1-\Phi\left(\frac c{\sigma/\sqrt n}\right)\right],
\]
بنابراین $c=z_{1-\alpha/2}\sigma/\sqrt n$.
\begin{example}{آزمون چپ‌طرفه}
داده‌ها عبارت‌اند از $n=16$، $\bar x=9$ و $\sigma^2=9$. فرض‌ها را
\[H_0:\mu=10,\qquad H_1:\mu<10\]
و سطح را $\alpha=0.05$ بگیرید. آنگاه
\[
z_{\rm obs}=\frac{9-10}{3/4}=-\frac43\approx-1.333.
\]
چون $-1.333>-1.64485$، فرض صفر رد نمی‌شود. مقدار $p$ حدود $0.09121$ است و همین تصمیم را تأیید می‌کند.
\end{example}
\begin{example}{آزمون دوطرفه با مقدار آمارهٔ ۲}
در مثال صفحهٔ ۳۱، $\bar x=3,\mu_0=1,\sigma^2=25$ و محاسبهٔ نوشته‌شده به $z_{\rm obs}=2$ می‌رسد. با خطای استاندارد $1$، حجم نمونهٔ متناظر $n=25$ است. برای
\[H_0:\mu=1,\qquad H_1:\mu\ne1,\qquad \alpha=0.05,\]
داریم $|2|>1.95996$؛ پس فرض صفر رد می‌شود. مقدار $p$ دوطرفه تقریباً $0.04550$ است.
\edit{زیرنویس حجم نمونه در صورت دست‌نویس کم‌خواناست؛ در این مثال، $n=25$ از خطای استاندارد و مقدار آمارهٔ همان حل بازسازی شده است.}
\end{example}
\section{آزمون میانگین یک جامعهٔ نرمال؛ واریانس نامعلوم}
\source{۳۲}
وقتی $\sigma$ نامعلوم باشد، آمارهٔ مناسب
\[T=\frac{\bar X-\mu_0}{S/\sqrt n}\sim t_{n-1}\quad\text{تحت }H_0\]
است. استخراج ناحیهٔ رد مانند حالت قبل انجام می‌شود و چندک‌های استیودنت جایگزین چندک‌های نرمال می‌شوند.
\begin{result}{سه آزمون استیودنت تک‌نمونه‌ای}
\[
\begin{array}{c|c}
H_1&\text{شرط رد }H_0\\\hline
\mu<\mu_0&t_{\rm obs}<-t_{1-\alpha,n-1}\\[2pt]
\mu>\mu_0&t_{\rm obs}>t_{1-\alpha,n-1}\\[2pt]
\mu\ne\mu_0&|t_{\rm obs}|>t_{1-\alpha/2,n-1}
\end{array}
\]
\end{result}
\begin{example}{آیا میانگین نمرات ریاضی برابر ۱۵ است؟}
در نمونه‌ای با حجم $9$، میانگین و انحراف معیار نمرات به‌ترتیب $13$ و $3$ است. با فرض نرمال بودن جامعه، برای
\[H_0:\mu=15,\qquad H_1:\mu\ne15\]
داریم
\[
t_{\rm obs}=\frac{13-15}{3/\sqrt9}=-2,\qquad
\nu=8,\qquad t_{0.975,8}=2.30600.
\]
چون $|-2|=2<2.30600$، در سطح ۵ درصد فرض صفر \textbf{رد نمی‌شود}. مقدار $p$ دوطرفه حدود $0.08052$ است.
\edit{در دست‌نویس مقدار بحرانی $1.86$، مربوط به $t_{0.95,8}$، برای آزمون دوطرفهٔ ۵ درصد استفاده و نتیجهٔ رد گرفته شده است. چندک درست $t_{0.975,8}$ و تصمیم درست، رد نکردن است.}
\end{example}
\section{آزمون اختلاف میانگین دو جامعهٔ مستقل}
\source{۳۳ و ۳۴}
در دست‌نویس، بعضی روابط با جهت $\mu_2-\mu_1$ نوشته شده‌اند. برای هماهنگی با فصل قبل، در این نسخه
\[\Delta=\mu_1-\mu_2,\qquad H_0:\Delta=\Delta_0\]
را به کار می‌بریم. تعویض جهت اختلاف، علامت آماره و جهت فرض‌های یک‌طرفه را هم عوض می‌کند.
\subsection{واریانس‌های معلوم}
آمارهٔ آزمون تحت فرض صفر برابر است با
\[
Z=\frac{(\bar X_1-\bar X_2)-\Delta_0}
{\sqrt{\sigma_1^2/n_1+\sigma_2^2/n_2}}\sim\N(0,1).
\]
برای $H_1:\Delta<\Delta_0$، $Z<-z_{1-\alpha}$؛ برای $H_1:\Delta>\Delta_0$، $Z>z_{1-\alpha}$؛ و برای $H_1:\Delta\ne\Delta_0$، $|Z|>z_{1-\alpha/2}$ شرط رد است.
\subsection{واریانس‌های نامعلوم ولی برابر}
از $S_p^2$ فصل قبل استفاده می‌کنیم:
\[
T_P=\frac{(\bar X_1-\bar X_2)-\Delta_0}
{S_p\sqrt{1/n_1+1/n_2}}\sim t_{n_1+n_2-2}.
\]
برای فرض مقابل کوچک‌تر، بزرگ‌تر و نابرابر، شرط‌های رد به‌ترتیب
\[
T_P<-t_{1-\alpha,\nu},\qquad
T_P>t_{1-\alpha,\nu},\qquad
|T_P|>t_{1-\alpha/2,\nu},\qquad \nu=n_1+n_2-2
\]
هستند.
\subsection{واریانس‌های نامعلوم و نامساوی}
\[
T_W=\frac{(\bar X_1-\bar X_2)-\Delta_0}
{\sqrt{S_1^2/n_1+S_2^2/n_2}}\approx t_\nu,
\]
که $\nu$ همان درجهٔ آزادی ساترویت در فصل قبل است. شرط‌های رد دقیقاً شکل سه شرط بالا را دارند، اما با $T_W$ و درجهٔ آزادی ولچ؛ کنترل سطح در این حالت تقریبی است.
\section{آزمون واریانس یک جامعهٔ نرمال}
\source{۳۵ و ۳۶}
فرض صفر $H_0:\sigma^2=\sigma_0^2$ است. بسته به معلوم بودن میانگین، دو آماره داریم:
\[
Q=\begin{cases}
\displaystyle\frac{\sum_i(X_i-\mu)^2}{\sigma_0^2},&\mu\text{ معلوم،}\quad\nu=n,\\[5pt]
\displaystyle\frac{(n-1)S^2}{\sigma_0^2},&\mu\text{ نامعلوم،}\quad\nu=n-1.
\end{cases}
\]
تحت فرض صفر، $Q\sim\chi^2_\nu$.

\textbf{استخراج برای فرض مقابل کوچک‌تر.} اگر $\mu$ معلوم باشد، قرار دهید $S_0^2=n^{-1}\sum_i(X_i-\mu)^2$. شرط رد را $S_0^2/\sigma_0^2<c$ بگیرید. از
\[
\alpha=\Prob_{H_0}\left(\frac{S_0^2}{\sigma_0^2}<c\right)
=\Prob(\chi^2_n<nc)
\]
نتیجه می‌شود $c=\chi^2_{\alpha,n}/n$. در حالت میانگین نامعلوم، جایگزینی $n$ با $n-1$ و $S_0^2$ با $S^2$ نتیجهٔ متناظر را می‌دهد.
\begin{result}{سه آزمون واریانس}
\[
\begin{array}{c|c}
H_1&\text{شرط رد }H_0\\\hline
\sigma^2<\sigma_0^2&Q<\chi^2_{\alpha,\nu}\\[2pt]
\sigma^2>\sigma_0^2&Q>\chi^2_{1-\alpha,\nu}\\[2pt]
\sigma^2\ne\sigma_0^2&Q<\chi^2_{\alpha/2,\nu}\ \text{یا}\ Q>\chi^2_{1-\alpha/2,\nu}
\end{array}
\]
در آزمون دوطرفهٔ این جزوه، احتمال هر دُم برابر $\alpha/2$ انتخاب شده است؛ چون کای‌دو نامتقارن است، شرط رد را با قدرمطلق یک حد متقارن نمی‌نویسیم.
\end{result}
\section{آزمون برابری واریانس دو جامعهٔ نرمال}
\source{۳۷ و ۳۸}
دو نمونهٔ نرمال مستقل در نظر بگیرید. از تعریف $R_j$ در فصل قبل استفاده می‌کنیم: اگر میانگین جامعه معلوم باشد، $R_j=n_j^{-1}\sum_i(X_{ji}-\mu_j)^2$ و $\nu_j=n_j$؛ در غیر این صورت $R_j=S_j^2$ و $\nu_j=n_j-1$.

برای فرض صفر $H_0:\sigma_1^2/\sigma_2^2=1$، آمارهٔ
\[F=\frac{R_1}{R_2}\sim F_{\nu_1,\nu_2}\]
تحت فرض صفر معتبر است. چون نسبت کوچک‌تر به سود واریانس اول کوچک‌تر است،
\begin{result}{سه آزمون نسبت واریانس‌ها}
\[
\begin{array}{c|c}
H_1&\text{شرط رد }H_0\\\hline
\sigma_1^2/\sigma_2^2<1&F<F_{\alpha;\nu_1,\nu_2}\\[2pt]
\sigma_1^2/\sigma_2^2>1&F>F_{1-\alpha;\nu_1,\nu_2}\\[2pt]
\sigma_1^2/\sigma_2^2\ne1&F<F_{\alpha/2;\nu_1,\nu_2}\ \text{یا}\ F>F_{1-\alpha/2;\nu_1,\nu_2}
\end{array}
\]
\end{result}
اگر فرض صفر $\sigma_1^2/\sigma_2^2=\rho_0$ باشد، آمارهٔ $F=(R_1/R_2)/\rho_0$ را با همین چندک‌ها مقایسه می‌کنیم.
\begin{important}{محدودیت مدل}
روابط دقیق کای‌دو و $F$ در این فصل به نرمال بودن جامعه‌ها متکی‌اند. صرف بزرگ بودن نمونه، فرض‌های دقیق این آزمون‌ها را جایگزین نمی‌کند. ترتیب صورت، مخرج و درجات آزادی را ثابت نگه دارید.
\end{important}
\section{آزمون میانگین در داده‌های زوجی}
\source{۳۹}
با تعریف $D_i=X_i-Y_i$، آزمون اختلاف میانگین‌ها به آزمون تک‌نمونه‌ای اختلاف‌ها تبدیل می‌شود. برای
\[H_0:\mu_D=\Delta_0\]
آمارهٔ
\[T_D=\frac{\bar D-\Delta_0}{S_D/\sqrt n}\sim t_{n-1}\]
را به کار می‌بریم؛ این رابطه در صورت نرمال بودن اختلاف‌ها دقیق است. در آزمون برابری میانگین‌ها $\Delta_0=0$ است. سه ناحیهٔ رد عبارت‌اند از
\[
\begin{array}{c|c}
H_1&\text{شرط رد }H_0\\\hline
\mu_D<\Delta_0&T_D<-t_{1-\alpha,n-1}\\[2pt]
\mu_D>\Delta_0&T_D>t_{1-\alpha,n-1}\\[2pt]
\mu_D\ne\Delta_0&|T_D|>t_{1-\alpha/2,n-1}
\end{array}
\]
استقلال لازم، استقلال \textbf{زوج‌ها از یکدیگر} است؛ استقلال دو عضو داخل هر زوج لازم نیست.
\section{آزمون نسبت یک جامعه}
\source{۴۰}
برای $X_i\iid\operatorname{Bern}(p)$ و $H_0:p=p_0$، با $0<p_0<1$، آمارهٔ آزمون بزرگ‌نمونه‌ای
\[Z=\frac{\widehat p-p_0}{\sqrt{p_0(1-p_0)/n}}\approx\N(0,1)\quad\text{تحت }H_0\]
است. برای فرض مقابل‌های $p<p_0$، $p>p_0$ و $p\ne p_0$ به‌ترتیب شرط‌های
\[
Z<-z_{1-\alpha},\qquad Z>z_{1-\alpha},\qquad |Z|>z_{1-\alpha/2}
\]
به کار می‌روند.
\begin{important}{در آزمون، واریانس را تحت فرض صفر بسازید}
برای این آزمون، خطای استاندارد شامل $p_0$ است؛ در فاصلهٔ والد از $\widehat p$ استفاده کردیم. شمار موفقیت و شکستِ مورد انتظار تحت فرض صفر باید برای تقریب نرمال کافی باشد. در نمونهٔ کوچک، آزمون دقیق دوجمله‌ای گزینهٔ مناسب‌تری است.
\end{important}
\section{آزمون برابری دو نسبت مستقل}
\source{۴۱ و ۴۲}
فرض کنید $X_{1i}\iid\operatorname{Bern}(p_1)$ و $X_{2j}\iid\operatorname{Bern}(p_2)$ و دو نمونه مستقل‌اند. هدف، آزمون
\[H_0:p_1=p_2\]
است. تحت فرض صفر، نسبت مشترک $p$ داریم و برآوردگر تجمیعی آن برابر است با
\[
\widehat p=\frac{\sum_iX_{1i}+\sum_jX_{2j}}{n_1+n_2}
=\frac{n_1\widehat p_1+n_2\widehat p_2}{n_1+n_2}.
\]
بنابراین آمارهٔ بزرگ‌نمونه‌ای آزمون برابر است با
\[
\boxed{Z=\frac{\widehat p_1-\widehat p_2}
{\sqrt{\widehat p(1-\widehat p)(1/n_1+1/n_2)}}.}
\]
برای فرض‌های مقابل $p_1<p_2$، $p_1>p_2$ و $p_1\ne p_2$، شرط‌های رد به‌ترتیب
\[
Z<-z_{1-\alpha},\qquad Z>z_{1-\alpha},\qquad |Z|>z_{1-\alpha/2}
\]
هستند. نسبت تجمیعی به این دلیل معتبر است که فرض صفر، برابری دو نسبت را تحمیل می‌کند. برای آزمون اختلاف غیرصفر $p_1-p_2=\Delta_0$، این تجمیع ساده عموماً معتبر نیست.
\section{مقدار احتمالی یا مقدار \lr{p}}
\source{۴۳ تا ۴۵}
\begin{definition}{مقدار احتمالی آزمون}
مقدار $p$ کوچک‌ترین سطح معنی‌داری در خانوادهٔ آزمون‌های مورد نظر است که در آن، دادهٔ مشاهده‌شده به رد $H_0$ منجر می‌شود. به‌طور معادل، تحت فرض صفر، احتمال مشاهدهٔ مقداری از آماره به اندازهٔ مقدار مشاهده‌شده یا افراطی‌تر در جهت فرض مقابل است.
\end{definition}
\begin{important}{مقدار $p$ چه چیزی را نمی‌گوید؟}
مقدار $p$ احتمال درست بودن $H_0$ نیست. این مقدار به مدل، فرض‌های آزمون و جهت فرض مقابل وابسته است؛ کوچک بودن آن نیز به‌تنهایی اندازه یا اهمیت عملی اثر را بیان نمی‌کند.
\end{important}
اگر آمارهٔ $T$ تحت فرض صفر تابع توزیع پیوستهٔ $G_0$ داشته باشد و مقدار مشاهده‌شده $t_{\rm obs}$ باشد،
\[
\begin{array}{c|c}
\text{جهت آزمون}&\text{مقدار }p\\\hline
\text{چپ‌طرفه}&G_0(t_{\rm obs})\\[2pt]
\text{راست‌طرفه}&1-G_0(t_{\rm obs})\\[2pt]
\text{دوطرفه با دو دُم برابر}&2\min\{G_0(t_{\rm obs}),1-G_0(t_{\rm obs})\}
\end{array}
\]
برای نرمال و استیودنتِ متقارن، مقدار دوطرفه را می‌توان به صورت $2[1-G_0(|t_{\rm obs}|)]$ نوشت. برای توزیع‌های نامتقارن، تعریف دوطرفه به شیوهٔ انتخاب ناحیهٔ رد وابسته است؛ فرمول جدول به آزمون دو دُم برابر مربوط می‌شود.
\begin{result}{قاعدهٔ تصمیم بر اساس مقدار $p$}
\[
p\le\alpha\quad\Longrightarrow\quad\text{رد }H_0,
\qquad
p>\alpha\quad\Longrightarrow\quad\text{رد نکردن }H_0.
\]
در مدل پیوسته، انتخاب $<$ یا $\le$ در مرز، احتمال خطا را تغییر نمی‌دهد؛ قرارداد تصمیم این جزوه $\le$ است.
\end{result}
\begin{example}{میانگین ضخامت محصول جدید}
ضخامت محصول قدیمی میانگین $500$ و انحراف معیار معلوم $20$ دارد. از محصول جدید، نمونه‌ای با $n=25$ و میانگین $508$ به دست آمده است. برای بررسی افزایش میانگین،
\[H_0:\mu=500,\qquad H_1:\mu>500,\qquad
z_{\rm obs}=\frac{508-500}{20/\sqrt{25}}=2.\]
با فرض مدل نرمال برای جامعه، یا مناسب بودن تقریب میانگین نمونه،
\[
p=\Prob(Z\ge2)=1-\Phi(2)=\boxed{0.0227501}.
\]
در سطح $\alpha=0.05$، چون $p<0.05$، فرض صفر رد می‌شود؛ در سطح $\alpha=0.01$، چون $p>0.01$، رد نمی‌شود.

همین نتیجه با مقادیر بحرانی به دست می‌آید: $z_{0.95}=1.64485<2$ ولی $z_{0.99}=2.32635>2$.
\end{example}
\begin{center}
\begin{LTR}
\begin{tikzpicture}[x=1.55cm,y=4cm]
\fill[teal!22] plot[domain=2:3.8,samples=65] (\x,{exp(-\x*\x/2)/sqrt(2*pi)})--(3.8,0)--(2,0)--cycle;
\draw[navy,thick] plot[domain=-3.4:3.8,samples=130] (\x,{exp(-\x*\x/2)/sqrt(2*pi)});
\draw[muted,-{Latex}] (-3.5,0)--(4,0) node[right] {$z$};
\draw[gold,dashed] (1.64485,0)--(1.64485,.29);
\draw[teal,thick] (2,0)--(2,.15);
\draw[muted,dashed] (2.32635,0)--(2.32635,.38);
\node[gold,anchor=south east] at (1.64485,.29) {$z_{0.95}=1.645$};
\node[muted,anchor=south west] at (2.32635,.38) {$z_{0.99}=2.326$};
\node[below,teal] at (2,0) {$z_{\rm obs}=2$};
\draw[-{Latex},teal] (3,.19) node[above] {$p=0.02275$}--(2.6,.012);
\end{tikzpicture}
\end{LTR}
\figcap{بازطراحی شکل‌های صفحه‌های ۴۳ تا ۴۵: مقدار مشاهده‌شده از حد بحرانی ۵ درصد عبور کرده، اما از حد بحرانی ۱ درصد عبور نکرده است. مساحت سبزآبی، مقدار $p$ را نشان می‌دهد.}
\end{center}
\section{رابطهٔ آزمون دوطرفه و فاصلهٔ اطمینان}
\begin{result}{دو بیان یک تصمیم}
در روش‌های متناظر، آزمون دوطرفهٔ سطح $\alpha$، فرض صفر $\theta=\theta_0$ را دقیقاً زمانی رد می‌کند که $\theta_0$ در فاصلهٔ اطمینان $1-\alpha$ متناظر نباشد. این هم‌ارزی به یکسان بودن مدل، آماره و روش ساخت دو طرف وابسته است.
\end{result}
برای مثال، در آزمون میانگین با واریانس معلوم،
\[
\left|\frac{\bar x-\mu_0}{\sigma/\sqrt n}\right|>z_{1-\alpha/2}
\quad\Longleftrightarrow\quad
\mu_0\notin\left[\bar x\pm z_{1-\alpha/2}\frac\sigma{\sqrt n}\right].
\]
\textbf{مسیر پیشنهادی حل مسائل آزمون.} ابتدا پارامتر و جهت اثر را مشخص کنید؛ فرض‌ها و سطح $\alpha$ را بنویسید؛ سپس شرایط روش و خطای استاندارد را تعیین کنید. آماره و مقدار $p$ یا حد بحرانی را محاسبه کنید و در پایان، تصمیم را در زبان مسئله همراه با جهت و اندازهٔ اثر گزارش دهید.



\end{document}
